\documentclass[10pt,twocolumn]{article}
\usepackage[T1]{fontenc}
\usepackage[utf8]{inputenc}
\usepackage{lmodern}
\usepackage[margin=1.8cm,top=2.4cm,bottom=2cm,columnsep=0.6cm]{geometry}
\usepackage{fancyhdr}
\usepackage{titlesec}
\usepackage{booktabs}
\usepackage{amsmath,amssymb,amsfonts}
\usepackage{microtype}
\usepackage{xcolor}
\usepackage{mdframed}
\usepackage{parskip}
\usepackage{enumitem}
\usepackage{hyperref}

\definecolor{rulered}{RGB}{180,30,30}
\definecolor{boxgray}{RGB}{245,245,245}
\definecolor{keywordgray}{RGB}{100,100,100}

\pagestyle{fancy}
\fancyhf{}
\fancyhead[L]{\textsc{\small Claude Mag}}
\fancyhead[C]{\small\textit{Machine Learning Research Digest}}
\fancyhead[R]{\small 2026-07-01}
\fancyfoot[C]{\small\thepage}
\renewcommand{\headrulewidth}{0.8pt}

\titleformat{\section}{\normalsize\bfseries\color{rulered}}{}{0pt}{}%
  [\vspace{-4pt}\color{rulered}\rule{\columnwidth}{0.4pt}]
\titlespacing{\section}{0pt}{8pt}{4pt}

\hypersetup{colorlinks=true,urlcolor=rulered,linkcolor=black}

\begin{document}
\setlength{\parindent}{0pt}
\setlength{\parskip}{4pt}

{\small\color{keywordgray}\textbf{Keywords:}
  interpretability \textbullet{}
  CLIP \textbullet{}
  concept bottleneck models \textbullet{}
  zero-shot classification}\par\vspace{4pt}

{\LARGE\bfseries\raggedright
  Multimodal Concept Bottleneck Models}\par\vspace{4pt}

{\large\itshape\raggedright
  Dual Concept Layers Make CLIP Fully Interpretable Without Sacrificing
  Accuracy---or Requiring Retraining for New Classes}\par\vspace{6pt}

{\small\textbf{By} Tongqing Shi, Ge Yan, Tuomas Oikarinen \& Tsui-Wei Weng
  (UC San Diego)}\par\vspace{2pt}

{\small\color{keywordgray}arXiv:2606.19882 \textbullet{} Submitted June 2026}
\par\vspace{2pt}\noindent{\color{rulered}\rule{\columnwidth}{0.6pt}}\vspace{6pt}

Interpretable image classifiers no longer need to sacrifice accuracy or be
retrained for every new class. MM-CBM introduces dual Concept Bottleneck
Layers for both the image and text branches of CLIP, forcing predictions to
pass through a human-readable concept space while matching black-box CLIP
accuracy within 5\%---and enabling zero-shot transfer to unseen classes
with no additional training.

\begin{mdframed}[backgroundcolor=boxgray,linecolor=rulered,linewidth=1pt,
    innerleftmargin=6pt,innerrightmargin=6pt,
    innertopmargin=6pt,innerbottommargin=6pt]
\textbf{\small AT A GLANCE}\par\vspace{4pt}
\begin{description}[leftmargin=0pt,labelsep=4pt,itemsep=2pt,parsep=0pt,
    font=\small\bfseries]
  \item[Problem:] Existing Concept Bottleneck Models are limited to fixed
    label sets and suffer from concept leakage, where predictive signals
    outside the intended concept space are inadvertently exploited.
  \item[Why it matters:] Interpretable multimodal models are critical for
    high-stakes domains (medical imaging, autonomous driving) where users
    need to understand and intervene on model reasoning.
  \item[Method:] Dual Concept Bottleneck Layers (CBLs) align both image
    and text CLIP embeddings into a shared interpretable concept space;
    classification is a pure concept-space dot product with no linear
    head, eliminating leakage by construction.
  \item[Result:] Up to 51.26\% accuracy improvement over prior CBMs;
    within 5\% of black-box CLIP; zero-shot transfer to unseen classes.
  \item[Evidence:] CUB-200-2011, ImageNet, Oxford Pets, Stanford Cars;
    vs.\ Label-free CBM, Post-hoc CBM, LaBo.
\end{description}
\end{mdframed}

\section{The Big Picture}

Deep neural networks are powerful but opaque. Concept Bottleneck Models
(CBMs) are a popular interpretability approach: they insert a bottleneck
layer whose activations correspond to named concepts (e.g.\ ``has wings,''
``metallic surface,'' ``striped''), making the model's reasoning transparent.
A user can inspect which concepts drove a prediction and directly intervene
by modifying concept scores.

Standard CBMs have two structural limitations. First, they are trained for
a fixed label space. If a new class appears, the model must be retrained.
Second, the final classification head (a linear layer over concept scores)
can exploit non-concept information encoded in concept activations---a
phenomenon called \emph{concept leakage}. The model appears interpretable
but silently uses signals invisible to the user.

MM-CBM addresses both problems by extending CBMs into CLIP's shared
vision-language embedding space. Two Concept Bottleneck Layers---one for
image embeddings, one for text embeddings---project both modalities into
a unified concept space. Classification is then a dot product between the
image concept vector and the text concept vector for each class description.
No separate classification head is needed; leakage is impossible because
the concept space fully mediates inference. New classes require only their
text description---no retraining.

\section{Why Existing Approaches Fall Short}

\textbf{Label-free CBM:} Uses CLIP to identify concept-image alignments
without manual concept labels, but still trains a linear head, preserving
leakage.

\textbf{Post-hoc CBM:} Applies concept bottleneck post-training by
projecting frozen representations onto concepts. Accuracy drops significantly
on fine-grained datasets.

\textbf{LaBo:} Uses language models to generate concepts and assembles
a bottleneck via logistic regression. Fixed to training classes; concept
leakage present via the regression head.

None of these methods operate purely in the concept space, and all
require retraining or re-optimisation for new classes.

\section{Core Mechanism}

Let $K$ concept descriptions $\{c_1,\ldots,c_K\}$ be given (e.g., from
a domain expert or generated by an LLM). Compute concept text embeddings
$e_k = f_T(c_k)$ using CLIP's text encoder $f_T$.

For an image $x_v$ with CLIP embedding $z_v = f_V(x_v)$, the image
concept vector is:
\[
[c_v]_k = \frac{\langle W_v\,z_v,\,e_k\rangle}
               {\|W_v\,z_v\|\cdot\|e_k\|},
\quad k=1,\ldots,K,
\]
where $W_v\in\mathbb{R}^{d\times d}$ is the learnable image CBL projection.
The text concept vector for class $l_j$ is:
\[
[c_t(l_j)]_k = \frac{\langle W_t\,z_t(l_j),\,e_k\rangle}
                    {\|W_t\,z_t(l_j)\|\cdot\|e_k\|}.
\]
Classification: $\hat y = \arg\max_j \langle c_v, c_t(l_j)\rangle$.
No linear head; no leakage.

\section{Technical Formulation}

\textbf{Training objective.}
Let $\mathcal{D}_\mathrm{aux}$ be an auxiliary dataset of image-text
pairs. The loss combines three terms:
\[
\mathcal{L} = \mathcal{L}_\mathrm{align}
            + \lambda_\mathrm{task}\,\mathcal{L}_\mathrm{task}
            + \lambda_\mathrm{sparse}\,\|c_v\|_1.
\]

$\mathcal{L}_\mathrm{align}$ is a CLIP-style contrastive loss on
$(c_v, c_t)$ pairs:
\[
\mathcal{L}_\mathrm{align} =
-\frac{1}{B}\sum_{b=1}^B\log
\frac{e^{\langle c_v^b,c_t^b\rangle/\tau}}
     {\sum_{b'} e^{\langle c_v^b,c_t^{b'}\rangle/\tau}}.
\]

$\mathcal{L}_\mathrm{task}$ is cross-entropy on training classes,
using $c_v$ as logits (after matching against class concept vectors).

$\|c_v\|_1$ is an $L_1$ sparsity penalty ensuring that each prediction
uses only a small number of active concepts.

\textbf{Zero-shot transfer.}
For a new class $l_\mathrm{new}$, compute $c_t(l_\mathrm{new})$
using the text encoder and CBL. No retraining required. This follows
from the universal concept space shared by image and text CBLs.

\textbf{User intervention.}
To intervene on concept $k$: set $[c_v]_k \leftarrow v$ for any
desired value $v \in [-1, 1]$. The corrected concept vector propagates
directly to the prediction via the dot product---fully transparent.

\section{Learning or Inference Procedure}

Training:
\begin{enumerate}[itemsep=2pt,parsep=0pt]
  \item Freeze CLIP encoders $f_V$, $f_T$.
  \item Initialise $W_v$, $W_t$ as identity.
  \item Train on $\mathcal{D}_\mathrm{aux}$ minimising $\mathcal{L}$
    with Adam optimiser ($5\times10^{-4}$ lr, 50 epochs).
\end{enumerate}
Inference: two forward passes ($f_V$ and $W_v$), one dot product per
class. Complexity $O(K \cdot d)$ per image---faster than CLIP zero-shot.

\section{What the Guarantee Says}

By construction, all predictive signal in MM-CBM flows through the
concept vectors $c_v$ and $c_t$. Since the final classifier is a dot
product in the concept space (no additional parameters), the model cannot
encode information outside the $K$-dimensional concept subspace. This is
a \textbf{zero-leakage} guarantee: any prediction not expressible via
concept scores is structurally impossible.

For zero-shot transfer, the guarantee is that any class expressible via
concept descriptions $\{c_k\}$ is immediately classifiable. Accuracy
on unseen classes depends on how well the concept vocabulary spans the
relevant visual dimensions---a design choice, not a training variable.

\section{Experimental Findings}

\begin{itemize}[itemsep=2pt,parsep=0pt]
  \item \textbf{CUB-200-2011:} MM-CBM achieves 78.3\% vs.\ 53.1\%
    for Label-free CBM (a 47.5\% relative gain) while staying within
    4\% of black-box CLIP.
  \item \textbf{ImageNet:} Within 5\% of CLIP zero-shot; best
    interpretable method.
  \item \textbf{Oxford Pets / Stanford Cars:} Consistent improvements
    over all CBM baselines (average 51.26\% accuracy improvement).
  \item \textbf{Intervention:} Correcting one concept score improves
    accuracy by 12\% on targeted error cases.
  \item \textbf{Zero-shot:} Correctly classifies 8 unseen bird species
    added at test time, with no accuracy degradation on original classes.
\end{itemize}

\section{Ablations and Interpretation}

\textbf{Remove image CBL} ($W_v = I$): Accuracy drops by $\approx8$
percentage points; concept alignment is worse.

\textbf{Remove text CBL} ($W_t = I$): Similar drop; text concept vectors
do not align with visual concept space without projection.

\textbf{Remove sparsity loss:} Accuracy $+0.3$\% but average concept
activations per prediction increase $8\times$, making interpretations
unintelligible.

\textbf{Concept set size $K$:} Performance plateaus at $K=100$;
smaller sets ($K<30$) significantly hurt accuracy on fine-grained tasks.

\section*{Reference}

Tongqing Shi, Ge Yan, Tuomas Oikarinen, Tsui-Wei Weng.
\textbf{Multimodal Concept Bottleneck Models}.
arXiv:2606.19882, June 2026.
\url{https://arxiv.org/abs/2606.19882}

\end{document}
